\section{Universal Algebra in \Frexlib{}}
\seclabel{frexlib}
To define an interface to algebraic simplifiers, we first
specify and represent algebraic structures.
We implement signatures and their operators in \Frexlib{} as follows
(below, left and middle):
\ThreeColDIY[]{.27}{
  \SignaturesPartA{}
  \SignaturesPartB{}
}{.35}{\OpDecl{}}
 {.35}{\MonoidOperationIntro{}}
The implementation uses \idris{}'s implicit record field for
\ArityName{}.  Users define concrete instances of \SignatureName{},
such as the signature \MonoidSig{} for monoids, by defining a
% injective
type family for the indexed field
\OpWithArityName{} (above, right).
%% Injectivity avoids projecting the arity in concrete cases, where
%% unification extracts it automatically.  Injectivity improves usability
%% but is not otherwise necessary.  Idris currently has limited
%% support for injectivity: unification makes use of the injectivity of
%% data-type constructors, but there is no means for a type to require a
%% judgementally injective type-level function.


\begin{figure*}
  \TwoColDIY{.375}{
      \Tuple{}
  \algebraOverDecl{}
  \AlgebraDecl{}
  }{.6}{
\CongruenceWRT{}
\SetoidAlgebra{}
  }
  %%
  %% and (c) their
  %% instantiation in the monoid frexlet
  %%\figlabel{monoid frexlet}
\caption{Algebras and setoid algebras in \Frexlib{}}
\figlabel{algebras}
\end{figure*}


\Frexlib{} represents the domain of an $n$-ary operation
with an $n$-ary vector (\figref{algebras}). (As in
\Haskell{},
backticks turn any name into an infix operator.)  For
example, the additive natural numbers form
an algebra for the monoid
signature as follows:
\MyCodeListing[0\baselineskip]{
  \PutListing{.5\textwidth}\NatAdditive{}
  }
Instead of specifying the \IdrisBound{Semantics} field directly,
this code uses the smart constructor \IdrisFunction{MkAlgebra},
which has a \IdrisBound{Sem} argument, instead of \IdrisBound{Semantics}.
This smart constructor transfers its \IdrisBound{Sem} argument into
\IdrisData{MakeAlgebra}'s \IdrisBound{Semantics} field by uncurrying
each \nvar{}-ary function into a function taking an \nvar{}-ary vector of arguments.
The departure from our usual naming scheme in which it is the constructor of
a record \IdrisType{R} that is called \IdrisData{MkR} indicates that the
smart constructor \IdrisFunction{MkAlgebra}, not the record constructor
\IdrisData{MakeAlgebra},
is the preferred way to construct
\IdrisType{Algebra} instances.
The \IdrisKeyword{\textbackslash{}case} keyword is an anonymous
function that immediately pattern-matches its argument.  \emph{Setoid
  algebras} further require an equivalence relation that forms a
congruence w.r.t.~the operations~(\figref{algebras}).


We implement terms over a signature as follows, mirroring their mathematical definition:
\MyCodeListing{
  \PutListing{.7\textwidth}\TermDecl}
Terms form an algebra, the \emph{free} algebra, with symbols denoting term formers:
\MyCodeListing{
\PutListing{.7\textwidth}\FreeDecl
}
Terms also form a monad, with \IdrisData{Done} as its unit and substitution
as its sequencing operation.

\begin{figure*}[t]
\MonoidAxiom{}

\MonoidTheory{}

\MonoidStructure{}

\MonoidModel{}
\caption{Axiomatising monoids in \Frexlib{}}
\figlabel{monoids}
\end{figure*}

Turning to equations, \Frexlib{} only needs
equations in a finite context, and we call its cardinality
the \emph{support} of the equation. We implement equations
and presentations as follows:
\ThreeColDIY{.26}{\EquationDecl[gobble=2]{}}{.3}{\PresentationDecl{}}
            {.38}{\AssociativityScheme[numbers=right,numbersep=-.1cm]{}}
\noindent\\[.3\baselineskip] For example, the monoid presentation
$\Monoid$ has three axioms: left and right neutrality, and
associativity. \Frexlib{} defines a generic collection of axiom
schemes (above, right).  The second argument to \IdrisFunction{EqSpec}
lists the arities \IdrisData{[n1, ..., nk]} of the operations in a
scheme, so on lines 1--2 the argument \IdrisData{[2]} in the type
\SchemeOps{} states that the scheme involves a single binary
operation.  The first argument \IdrisData{3} in \MkEquation{} (lines 5--6)
indicates that the declaration of the scheme involves three variables
\IdrisFunction{X}\;\IdrisData{0}, \IdrisFunction{X}\;\IdrisData{1}, and
\IdrisFunction{X}\;\IdrisData{2}. We use these
generic axiom schemes to construct, for example, the theory of monoids in the
middle of \figref{monoids}.


\begin{figure*}[t]
\Models{}

\medskip

\Entails{}

\medskip

\ValidatesEquation{}

\medskip

\Validates{}
  \caption{Equational validity in an algebra}
  \figlabel{equational validity}
\end{figure*}

\figref{equational validity} shows
\Frexlib{}'s representation of what it means for an algebra
to validate an equation. We use \idris{}'s dependent
pairing construct to pair an algebra with an environment in the
standard entailment syntax \ModelEnv{}.
The following code validates the monoid axioms
for our running example:
\MyCodeListing{
  \PutListing{.5\textwidth}\AdditiveNatValid{}
}
We define models for a presentation:
\MyCodeListing{
  \PutListing{.5\textwidth}\ModelDecl{}
}
%
We can now define a monoid to be a $\Monoid$-model, as in
\figref{monoids}. For another example, now putting everything together,
we validate the monoid structure of multiplication as follows:
\MyCodeListing{
  \PutListing{.55\textwidth}{\NatMultitive[numbers=right]}
}
\label{cast explanation}
The coercion function %
(\IdrisFunction{cast} \IdrisKeyword{:} \IdrisBound{from} \IdrisKeyword{->} \IdrisBound{to}) %
on Line 3 is a method in the standard Idris type-class/interface
\IdrisType{Cast}
\IdrisBound{from} \IdrisBound{to}.  \Frexlib{} exports a
\IdrisType{Cast}-instance that converts an algebra into a setoid
algebra whose equivalence relation is propositional equality
(\IdrisKeyword{=}).  Lines 7--9 use results about the natural numbers
from \idris{}'s standard library.

\begin{figure}
  \ListInvMonoid{}
  \caption{The involutive monoids of list reversal}
  \figlabel{involutive lists}
\end{figure}

\paragraph{Using \Frexlib{}.}
Unless they are already working abstractly with an algebraic structure,
we expect that in practice users start by recognising that their concrete
algebra validates the axioms of an existing simplification
module---\emph{frexlet} for short. Such modules export a presentation,
convenient notation suites for its signature, and fral and/or frex simplifiers
for this presentation.

As a concrete example, we will take computations with lists
that also involve the \IdrisFunction{reverse} function. These form an
\emph{involutive} monoid: a monoid $\A$ equipped with a unary
\emph{involution} operator $x\mapsto\overline x : \U\A \to \U\A$
satisfying two axioms $\overline{\overline \x} = \x$ and
$\overline{\x\y} = \overline\y\,\overline\x$. We then equip our type
of interest, lists, with an involutive structure as in
\figref{involutive lists}. We can use this algebra and the involutive
monoid to discharge equations containing list variables and concrete
lists: \InvMonoidFrexExample[numbers=right, numbersep=-.125cm]

The \IdrisFunction{solve} function takes as argument the number of
variables (\IdrisBound{n}=\IdrisData{2} on line~2) in the algebraic
term to simplify, and an algebraic simplifier from the frexlet
(\IdrisFunction{Involutive.Frex.Frex} on line~4). The final argument
is a pair of terms with \IdrisBound{n}=\IdrisData{2} variables
(\IdrisFunction{Dyn} \IdrisData{0} and \IdrisFunction{Dyn}
\IdrisData{1}) and concrete values from the algebra. By importing
notation modules the frexlet provides, we can use infix multiplicative
notation such as (\IdrisFunction{.*.}). The type-checker then infers
the terms to substitute for each variable.

In this example, we used \IdrisFunction{solve} to define a stand-alone
lemma, but we may also call \IdrisFunction{solve} directly from a
chain of equational reasoning steps. When we extract lemmas, we often
want to prove them more abstractly, for \emph{all} involutive
monoids. In that case we use a fral:
\InvMonoidFralExample[numbers=right,numbersep=-.125cm] Lines 2--3 and 8--9 overload the
infix and postfix notation using the frexlet's built-in notation
suites. Concretely, the projection \IdrisFunction{Notation1} brings
into scope the functions
\IdrisKeyword{(}\IdrisFunction{.*.}\IdrisKeyword{)} and
\IdrisKeyword{(}\IdrisFunction{.inv}\IdrisKeyword{)} when writing
algebraic terms.  The \IdrisFunction{solve} function takes the number
of free variables and a corresponding fral simplifier (line~10), as
well as the two terms representing the equation of interest.  The
variables \IdrisBound{x}, \IdrisBound{y}, \IdrisBound{z} (bound in line 7) are
implicitly used in this call. \subsecref{solver implementation}
covers the type of \IdrisFunction{solve} in more detail.
